\documentclass[11pt]{article}
\usepackage[margin=1in]{geometry}
\usepackage{natbib}
\usepackage{booktabs}
\usepackage{amsmath}
\usepackage{graphicx}
\usepackage{hyperref}
\usepackage{lineno}

\title{Time-, tissue- and treatment-associated heterogeneity in tumour-residing migratory dendritic cells}
\author{
G\'{a}lvez-Cancino, F.$^{1,\dagger}$ \and Simpson, A.P.$^{1,\dagger}$ \and Lees, K.A.$^{1}$ \and Lisko, C.$^{1}$ \and Bawa, P.$^{1}$ \and Guerin, M.V.$^{1}$ \and Salgado, R.$^{2}$ \and Quezada, S.A.$^{1,*}$\\[6pt]
\small $^{1}$Cancer Immunology Unit, Research Department of Haematology, UCL Cancer Institute, University College London, London, UK\\
\small $^{2}$Department of Pathology, GZA-ZNA Hospitals, Antwerp, Belgium\\
\small $^{\dagger}$These authors contributed equally\\
\small $^{*}$Corresponding author
}
\date{}

\begin{document}
\maketitle

\begin{abstract}
Tumour dendritic cells (DCs) internalise antigen and upregulate CCR7, which directs their migration to tumour-draining lymph nodes (dLN). CCR7 expression is coupled to a maturation programme enriched in regulatory molecule expression, including PD-L1, termed mRegDC. However, the spatio-temporal dynamics and role of mRegDCs in anti-tumour immune responses remain unclear. Using photoconvertible mice to precisely track DC migration, we found that mRegDCs were the dominant DC population arriving in the dLN, but a subset remained tumour-resident despite CCR7 expression. These tumour-retained mRegDCs were phenotypically and transcriptionally distinct from their dLN counterparts and were heterogeneous. Specifically, they demonstrated a progressive reduction in the expression of antigen presentation and pro-inflammatory transcripts with more prolonged tumour dwell-time. Tumour mRegDCs spatially co-localised with PD-1$^+$CD8$^+$ T cells in human and murine solid tumours. Following anti-PD-L1 treatment, tumour-residing mRegDCs adopted a state enriched in lymphocyte stimulatory molecules, including OX40L, which was capable of augmenting anti-tumour cytolytic activity. Altogether, these data uncover previously unappreciated heterogeneity in mRegDCs that may underpin a variable capacity to support intratumoural cytotoxic T cells, and provide insights into their role in cancer immunotherapy.
\end{abstract}

\section{Introduction}

Dendritic cells (DCs) capture tumour antigens and upregulate the chemokine receptor CCR7, which directs their migration to secondary lymphoid organs where they present antigens to T cells \citep{sallusto1998, roberts2016}. Two major conventional DC (cDC) subsets have been identified in tumours: cDC1 that specialise in cross-presenting tumour antigens to CD8$^+$ T cells and are associated with improved survival \citep{bottcher2018}, and cDC2 that primarily activate CD4$^+$ T cells. The recent application of single-cell technologies to tissue samples has enabled the distinction of an intratumoral DC population that may arise from both cDC1 and cDC2, but demonstrates a conserved programme characterised by high expression of CCR7, the costimulatory molecules CD80 and CD86, and importantly, the immunoregulatory molecules PD-L1 and PD-L2 \citep{maier2020, gerhard2021}.

This DC subset has been variously labelled as migratory DCs, mature DCs enriched in immunoregulatory molecules (mRegDCs, the term we adopted in the current study), as well as LAMP3$^+$ DCs \citep{nieto2022}. Regardless of nomenclature, single-cell atlasing efforts have identified mRegDCs in multiple human cancers \citep{gerhard2021, cheng2021, zilionis2019}, and the acquisition of this maturation programme appears dependent on the uptake of tumour antigen \citep{maier2020}. The presence of mature LAMP$^+$ DCs within lymphoid follicles has been associated with improved prognosis in non-small cell lung cancer (NSCLC) \citep{dieu2004, dieu2008, ladanyi2014}. However, the spatio-temporal dynamics and the functional role of mRegDCs in anti-tumour immunity remain poorly understood.

Critically, PD-L1 expression by DCs has been shown to be essential for effective responses to anti-PD-L1 antibodies \citep{peng2020, oh2020}, suggesting that mRegDCs may be important mediators of immune checkpoint blockade (ICB) therapy. Yet, whether mRegDCs are functionally homogeneous, and whether their role changes with duration in the tumour microenvironment (TME), has not been investigated.

In this study, we use photoconvertible Kaede transgenic mice to precisely track DC migration in tumour-bearing hosts \citep{tomura2008, galvez2023}. We discover that despite expressing CCR7, a substantial proportion of mRegDCs remain resident in the tumour, and these cells are transcriptionally distinct from their dLN counterparts. We uncover a time-dependent maturation trajectory in which tumour-retained mRegDCs progressively lose antigen presentation capacity and acquire features reminiscent of T cell exhaustion \citep{wherry2011}. Importantly, anti-PD-L1 treatment drives a subset of tumour mRegDCs toward an activated state enriched in lymphocyte costimulatory molecules, including OX40L, 4-1BBL, CD70, and PVR, which enhances anti-tumour CD8$^+$ T cell responses.

\section{Results}

\subsection{mRegDC signatures are associated with improved survival in human cancers}

The effect of tumour mRegDCs on prognosis in human cancers has not been broadly examined. To address this, we analysed the transcriptomes of 4,045 human solid tumours from The Cancer Genome Atlas (TCGA) \citep{cancer2023}. Kaplan-Meier analysis of overall survival stratified by enrichment of mRegDC signature genes revealed associations with improved survival in skin cutaneous melanoma (SKCM), colorectal adenocarcinoma (COAD), breast invasive carcinoma (BRCA), and lung adenocarcinoma (LUAD), all of which harbour CCR7$^+$ mRegDCs. Further analysis of 1,853 human breast tumours from METABRIC \citep{curtis2012} revealed cancer subtype-specific associations with survival. These data suggest that mRegDCs contribute to anti-tumour responses across a range of human cancers.

\subsection{mRegDCs are transcriptionally heterogeneous, and some remain within the tumour despite CCR7 expression}

To investigate the mechanisms by which mRegDCs promote anti-tumour immunity, we established multiple syngeneic subcutaneous colorectal tumours (MC38, MC38-Ova, CT26) in photoconvertible Kaede transgenic mice \citep{tomura2008}. Tumours were transcutaneously photoconverted on day 13, converting all infiltrated cells from the default green fluorescence (Kaede-green) to a red fluorescent profile (Kaede-red). At 24--72 hours following photoconversion, tumours were harvested, enabling the distinction of newly-infiltrating Kaede-green cells and Kaede-red cells retained in the tumour.

Single-cell RNA sequencing (scRNA-seq) of FACS-sorted CD45$^+$ cells generated 80,556 single-cell transcriptomes following rigorous quality control, including 32,191 myeloid cells. Unbiased clustering of DCs revealed 8 distinct clusters, including cycling DCs, cDC1, cDC2s, and mRegDCs, with mRegDC showing high expression of \textit{Ccr7}, \textit{Cd274} and \textit{Pdcd1lg2}. CCR7$^+$PD-L2$^+$ mRegDCs were a prominent tumour DC population. Surprisingly, tumour mRegDCs were predominantly Kaede-red, indicating that they have resided in the tumour for over 48 hours. In contrast, cDC1 and cDC2 were mostly Kaede-green, consistent with newly infiltrating cDCs differentiating into mRegDCs following uptake of tumour antigen \citep{maier2020}.

Pseudotime analysis rooted in the Kaede-green dominant cluster (cDC2\_1) revealed a trajectory terminating in mRegDC\_2/3 that transitioned through an intermediate mRegDC\_1 state, consistent with the Kaede-green/red ratio of clusters. RNA velocity analysis confirmed the cell-state transition from cDC, through mRegDC\_1, to mRegDC\_2 or mRegDC\_3. The receptor tyrosine kinase Axl, which recognises apoptotic cells and induces the mRegDC programme \citep{maier2020}, was upregulated as cDC2\_1 transitioned to cDC2\_3.

\subsection{Migrated mRegDCs are phenotypically distinct from tumour-resident populations}

Kaede-red CCR7$^+$PD-L2$^+$ DCs were readily detectable in the dLN, but not in the contralateral non-draining LN. Among Kaede-red DCs in the dLN, essentially all were mRegDC, up to 72 hours post-photoconversion. scRNA-seq of Kaede-red CD45$^+$ cells in tumour-dLNs showed that 94\% of Kaede-red myeloid cells in the dLN were mRegDCs. Hence, mRegDCs are the dominant myeloid cell tumour emigrants arriving in the dLN.

Gene set enrichment analysis (GSEA) and differential gene expression analysis revealed that tumour-residing mRegDCs and dLN mRegDCs were transcriptionally distinct. CD80 and CD86 were higher in tumour-residing mRegDCs, while \textit{Icosl} was significantly higher in migrated mRegDCs in the dLN. \textit{Il12b}, which drives anti-tumour cytotoxic T cell activity, was enriched in tumour mRegDCs, while \textit{Il15ra} was higher in mRegDC in the dLN \citep{riley2002}.

Strikingly, over 80\% of Kaede-red DCs from the dLN mapped to the tumour mRegDC\_1 cluster. Only 9\% of DCs in the dLN resembled mRegDC\_2/3, despite mRegDC\_2/3 being the dominant Kaede-red tumour mRegDC states. We propose that mRegDC\_2 and mRegDC\_3 are terminal, tumour-residing mRegDC states, and DCs that have transitioned beyond the intermediate mRegDC\_1 state become increasingly unlikely to egress.

\subsection{Tumour-retained mRegDCs progress towards an exhausted-like state, attenuated by anti-PD-L1 treatment}

During the progression from mRegDC\_1 to mRegDC\_3, there was a decrease in MHC-II expression, including class-II invariant chain \textit{Cd74}. Concomitantly, DC antigen presentation genes decreased as cells transitioned towards the tumour-retained mRegDC\_2/3 states, including downregulation of molecules involved in antigen processing (\textit{Ctsb}, \textit{Ctsl}, \textit{Ifi30}, \textit{Lgmn}, \textit{Psme2}), transport (\textit{Tap1}, \textit{Tap2}), chaperones (\textit{Hsp1a1}, \textit{Hspa2}, \textit{Hspa8}), MHC loading (\textit{H2-DMa}, \textit{H2-Oa}, \textit{Tapbp}, \textit{Calr}) and \textit{Ciita} \citep{chang1995}.

Moreover, genes involved in innate immune function and response to inflammatory stimuli, including interferon gamma/alpha response, inflammatory response, and Toll-like receptor signalling pathways, significantly decreased as cells transitioned from mRegDC\_1 to mRegDC\_3. Specifically, mRegDC\_3 showed reduced expression of several innate immune response genes (\textit{Nlrp3}, \textit{Tlr1}, \textit{Cd14}), genes involved in DC activation and migration (\textit{Axl}, \textit{Ccr7}, \textit{Rgs1}, \textit{Icam1}), lymphocyte co-stimulation (\textit{Cd40}, \textit{Tnfsf9}, \textit{Pvr}) and cytokines (\textit{Cxcl9}, \textit{Il1b}).

Anti-PD-L1 treatment influenced tumour mRegDC heterogeneity significantly. Assessment of differential abundance showed significant enrichment within mRegDC\_2 neighbourhoods following anti-PD-L1 treatment, but relative depletion of mRegDC\_3. The transcriptome of mRegDC\_2 was highly activated compared to mRegDC\_3, with increased expression of many genes involved in immune signalling, including \textit{Cd40}, \textit{Il1b}, \textit{Tnf}, and \textit{Nfkbia}. Several lymphocyte activation ligands \citep{croft2009, johnston2015} were significantly upregulated on mRegDC\_2 versus other tumour mRegDC states, such as \textit{Il12b}, \textit{Tnfsf4} (OX40L), \textit{Tnfsf9} (4-1BBL), \textit{Cd70} and \textit{Pvr}.

\subsection{mRegDCs interact with CD8$^+$ T cells in human tumours}

We deconvolved 521 human CRC transcriptomes and found the highest correlation between mRegDC and CD8$^+$ T cell transcripts \citep{becht2016}. CellPhoneDB analysis \citep{efremova2020} of scRNA-seq of human CRC \citep{pelka2021} showed that the CCR7$^+$CD274$^+$PDCD1LG2$^+$ DC cluster had stronger predicted interactions with effector CD8$^+$ T cells than other myeloid cells, including via CD28, CTLA-4 and PD-1 engagement.

Spatial transcriptomics data from human CRC tumours revealed hotspots with co-localised expression of mRegDC and effector CD8$^+$ T cell genes. Spatial correlation of mRegDC and effector CD8$^+$ T cell transcripts was also evident in human melanoma and breast tumours. Analysis of PIC-seq data from NSCLC \citep{tietscher2023} confirmed that mRegDC and effector CD8$^+$ T cells physically interact, with mRegDC-CD8$^+$ T cell doublets more frequent than other mRegDC-T cell combinations.

\subsection{Anti-PD-L1 promotes immunogenic mRegDC-CD8$^+$ T cell interactions}

Cell-cell communication analysis showed several interaction pairs including TNFSF4--TNFRSF4 (OX40L--OX40), TNFSF9--TNFRSF9 (4-1BBL--4-1BB), CD70--CD27, and PVR-mediated interactions upregulated between mRegDC\_2 and CD8$^+$ T$_{EX}$ cells. Functional studies using bone marrow-derived mRegDC-like cells co-cultured with OT-I CD8$^+$ T cells demonstrated that anti-PD-L1 treatment increased activation, clonal expansion, and granzyme B production. OX40L-deficient mRegDC-like cells showed reduced capacity to activate OT-I cells, confirming the importance of the OX40L:OX40 axis.

IFN$\gamma$ treatment of mRegDC-like cells upregulated expression of OX40L, PVR and CD40, consistent with the mRegDC\_2 state, and increased the activation and expansion of OT-I cells. This effect was reduced in cultures containing OX40L-deficient cells.

\subsection{Conserved mRegDC heterogeneity and CD8$^+$ T cell crosstalk in human cancers}

In scRNA-seq of human CRC, mRegDCs showed a gradient of MHC-II expression. Importantly, markers upregulated on tumour-residing MHC-II$^{\text{low}}$ mRegDCs in murine tumours were also preferentially expressed in human MHC-II$^{\text{low}}$ mRegDCs, including IL15, PVR, TNFSF4, TNFSF9, and CD70.

Analysis of 208 bulk transcriptomes from the IMvigor210 atezolizumab trial for metastatic urothelial carcinoma \citep{mariathasan2018} using Scissor integration \citep{sun2021} identified myeloid cells associated with a favourable clinical response that expressed transcripts associated with tumour-residing mRegDCs, including \textit{CCR7}, \textit{CD274}, \textit{IL15}, \textit{PVR} and \textit{CD70}. In scRNA-seq of breast cancer \citep{bassez2021}, mRegDCs from tumours that successfully underwent T cell clonotype expansion following anti-PD-1 treatment were enriched in transcripts associated with the ICB-induced mRegDC\_2 cluster.

\section{Discussion}

Our study reveals previously unappreciated heterogeneity in mRegDCs that has important implications for understanding anti-tumour immunity and the mechanisms of immune checkpoint blockade therapy. Using photoconvertible Kaede mice to track DC migration in real-time, we demonstrate that despite expressing CCR7, a substantial proportion of mRegDCs remain in the tumour microenvironment, where they undergo time-dependent transcriptional changes.

The key finding is that tumour-retained mRegDCs are not static; they progress through distinct transcriptional states. Newly formed mRegDC\_1 cells retain antigen presentation capacity and are the primary population that migrates to the dLN. In contrast, mRegDC\_2 and mRegDC\_3 represent terminal tumour-resident states with progressively diminished antigen presentation and inflammatory signalling---features reminiscent of T cell exhaustion \citep{wherry2011}. The concept of DC exhaustion extends recent observations of similar phenomena in NK cells within the TME \citep{galvez2023}.

Critically, anti-PD-L1 treatment redirected the mRegDC trajectory toward the mRegDC\_2 activated state rather than the exhausted mRegDC\_3 state. This mRegDC\_2 state was characterised by upregulation of lymphocyte stimulatory molecules, particularly OX40L, 4-1BBL, CD70, and PVR \citep{croft2009}. Our functional studies confirmed that OX40L expression on mRegDCs is important for augmenting anti-tumour CD8$^+$ T cell responses, and that IFN$\gamma$ from activated T cells may create a positive feedback loop that sustains the activated mRegDC state.

The conservation of mRegDC heterogeneity and mRegDC-CD8$^+$ T cell crosstalk across murine and human solid tumours, including CRC, melanoma, breast cancer, NSCLC, and urothelial carcinoma, supports the broad relevance of these findings. The association between the activated mRegDC state and clinical responses to ICB therapy suggests that tumour-residing mRegDCs may serve as both biomarkers and therapeutic targets for improving immunotherapy outcomes.

\section{Methods}

\subsection{Mice and tumour models}
Transgenic C57BL/6 Kaede, BALB/c Kaede, OX40L$^{+/\text{Human-CD4}}$ reporter, OX40L$^{\text{fl/fl}}$, and CD11c$^{\text{cre}}$ OX40L$^{\text{fl/fl}}$ mice were maintained at the University of Birmingham Biomedical Services Unit. MC38, MC38-Ova, and CT26 murine colon adenocarcinoma cells ($2.5 \times 10^5$) were subcutaneously injected in 100~$\mu$l PBS. Anti-PD-L1 antibodies were administered on days 7, 10 and 13.

\subsection{Photoconversion and cell tracking}
Tumours were transcutaneously photoconverted on day 13, converting Kaede-green to Kaede-red fluorescence. Tumours and dLNs were harvested 5h, 24h, 48h, or 72h after photoconversion.

\subsection{Single-cell RNA sequencing}
FACS-sorted CD45$^+$ cells were processed for scRNA-seq. Myeloid cells were clustered and annotated based on canonical marker genes and published signatures \citep{maier2020, gerhard2021}. Pseudotime and RNA velocity analyses were performed. Differential gene expression, GSEA, and CellPhoneDB analyses were conducted to characterise mRegDC heterogeneity and cell-cell interactions.

\subsection{Human tumour analyses}
Bulk transcriptomes from TCGA ($n = 4{,}045$), METABRIC ($n = 1{,}853$), and IMvigor210 ($n = 208$) were analysed. scRNA-seq datasets from human CRC \citep{pelka2021}, breast cancer \citep{bassez2021}, and urothelial carcinoma were integrated. Spatial transcriptomics (10X Genomics Visium) of CRC, melanoma, and breast tumours were used to assess mRegDC-CD8$^+$ T cell co-localisation. PIC-seq data from NSCLC were analysed for myeloid-T cell doublets.

\subsection{In vitro and ex vivo functional assays}
Bone marrow-derived DCs were cultured with UV-irradiated MC38-Ova cells to generate mRegDC-like cells. These were co-cultured with naive OT-I CD8$^+$ T cells for 3 days with or without anti-PD-L1 antibodies. OX40L-deficient mRegDC-like cells were generated from CD11c$^{\text{cre}}$ Tnfsf4$^{\text{fl/fl}}$ mice. Recombinant IFN$\gamma$ was used to assess DC-extrinsic effects.

\bibliographystyle{plainnat}
\bibliography{references}

\end{document}
